If $a\in S_{\mathfrak{p}}(0)$, then $sa=0$ for some $s\notin\mathfrak{p}$. In any $\mathfrak{p}$-primary ideal $\mathfrak{q}$, primaryness therefore gives $a\in\mathfrak{q}$.

Proceed by induction on $n$, taking $\mathfrak{p}_1$ itself when $n=1$. Reorder so that $\mathfrak{p}_n$ is maximal among the given primes, and take a minimal primary decomposition $\mathfrak{b}=\bigcap_{i<n}\mathfrak{q}_i$ with $r(\mathfrak{q}_i)=\mathfrak{p}_i$. If $\mathfrak{b}\subseteq S_{\mathfrak{p}_n}(0)$, choose a minimal prime $\mathfrak{p}\subseteq\mathfrak{p}_n$. By \exref{4.10}, $\mathfrak{b}\subseteq S_{\mathfrak{p}}(0)$ and therefore $\bigcap_{i<n}\mathfrak{p}_i\subseteq\mathfrak{p}$. Some $\mathfrak{p}_i$ then lies in $\mathfrak{p}$, forcing $\mathfrak{p}_i=\mathfrak{p}$, a contradiction.

Hence $\mathfrak{b}\not\subseteq S_{\mathfrak{p}_n}(0)$. By the assumed intersection property, choose a $\mathfrak{p}_n$-primary ideal $\mathfrak{q}_n$ with $\mathfrak{b}\not\subseteq\mathfrak{q}_n$. Set $\mathfrak{a}=\bigcap_{i=1}^n\mathfrak{q}_i$. The last component is indispensable. For $i<n$, choose $x\in\bigcap_{j<n,\,j\ne i}\mathfrak{q}_j\setminus\mathfrak{q}_i$ and $y\in\mathfrak{p}_n\setminus\mathfrak{p}_i$. For a sufficiently large $m$, $xy^m$ belongs to every $\mathfrak{q}_j$ with $j\ne i$, but not to $\mathfrak{q}_i$. Thus the decomposition is minimal and has precisely the required associated primes.
